% HW3 architecture diagram, one A4 landscape page. Build with ./build.sh (converts the SVG first).
\documentclass{article}
\usepackage[a4paper,landscape,margin=12mm,top=11mm,bottom=10mm]{geometry}
\usepackage{fontspec}
\usepackage{xeCJK}
\usepackage{graphicx}
\usepackage{xcolor}
\usepackage{enumitem}

\setmainfont{SF Pro Text}[BoldFont={SF Pro Text Semibold}]
\setmonofont{Menlo}[Scale=0.9]
% PingFang.ttc holds HK and TC faces; by name XeTeX picks the first (HK), so select TC by index.
\setCJKmainfont{PingFang.ttc}[FontIndex=2, BoldFont=PingFang.ttc, BoldFeatures={FontIndex=10}]
\newfontfamily\titlefont{SF Pro Display Semibold}
\newCJKfontfamily\titlecjk{PingFang.ttc}[FontIndex=10]
\xeCJKsetup{CJKecglue={\hskip 0.25em plus 0.1em}}
\xeCJKDeclareCharClass{Default}{"2013,"2022,"2212,"2192,"00D7}

\definecolor{muted}{HTML}{5B6575}
\setlength{\parindent}{0pt}
\AtBeginDocument{\fontsize{10.5}{16}\selectfont}
\pagestyle{empty}

\begin{document}

{\titlefont\titlecjk\fontsize{17}{21}\selectfont Homework \#3 系統架構}\hfill
{\color{muted}\small 黃毓峰\quad 315834011\quad Cloud Computing \& Services Architecture}

\vspace{1.5mm}
{\color{muted}\small 兩節點 K3s 叢集上的三層式網頁應用：一台 control-plane、一台 worker，web、api、db 各一個 replica，PostgreSQL 資料放在 PVC 上。}

\vspace{2mm}
\centering
\includegraphics[width=0.9\linewidth]{architecture-diagram.pdf}

\vspace{2mm}
\raggedright
\begin{itemize}[leftmargin=1.2em,itemsep=0.4mm,topsep=0mm]
  \item \textbf{請求路徑}\hspace{0.6em}瀏覽器連 web 的 NodePort 30080，nginx 把 \texttt{/api/} 轉給 Service \texttt{api}，Flask 經 Service \texttt{db} 連 PostgreSQL，各層只寫 Service 名稱。
  \item \textbf{持久儲存}\hspace{0.6em}db-0 掛載 PVC（PersistentVolumeClaim，儲存空間申請）\texttt{data-db-0}，資料在 worker 磁碟上，Pod 重建或關機重開後仍在。
  \item \textbf{節點分工}\hspace{0.6em}hw3-server 只跑 K3s 控制元件，hw3-worker 執行三層 App。兩台是 Incus 建立的 KVM 虛擬機，以 \texttt{nycu-cc-hw3-br} 互連。
\end{itemize}

\end{document}
